% META: {"id": "TSPE-CONTIN-CO-042", "chapitre": "TSPE-CONTINUITE", "type_objet": "corrige", "exercice_ref": "TSPE-CONTIN-EX-042", "capacites_codes": ["C2"], "sources_inspiration": [], "status": "generated"}
% BEGIN-VERIFY
% from sympy import symbols, diff, Rational, sqrt, simplify, solveset, S as SS, Eq, nsolve, Abs
% x = symbols('x')
% def iterer(f, u0, n):
%     u = u0
%     out = [u]
%     for _ in range(n):
%         u = simplify(f.subs(x, u))
%         out.append(u)
%     return out
% f = Rational(3,4)*x + 5
% assert solveset(Eq(f, x), x, SS.Reals) == {20}
% t = iterer(f, 40, 4)
% assert [float(v) for v in t] == [40.0, 35.0, 31.25, 28.4375, 26.328125]
% assert all(20 <= float(f.subs(x, v)) <= 40 for v in (20, 30, 40))
% assert simplify((f - x) + (x - 20)*Rational(1,4)) == 0
% for k in range(5):
%     assert abs(float(t[k]) - (20 + 20*0.75**k)) < 1e-12
% n = 0
% while 20 + 20*0.75**n >= 21:
%     n += 1
% assert n == 11
% assert 20 + 20*0.75**10 >= 21 > 20 + 20*0.75**11
% END-VERIFY
\begin{corrige}{TSPE-CONTIN-EX-042}

\textbf{Q1.} On conserve les trois quarts de la concentration precedente et on ajoute un quart d'une solution a $20$ g$\cdot$L$^{-1}$ : $s_{n+1} = 0{,}75\,s_n + 0{,}25 \times 20 = 0{,}75\,s_n + 5$.

$s_1 = 35$ ; $s_2 = 31{,}25$ ; $s_3 = 28{,}4375$ ; $s_4 = 26{,}328125$.

\textbf{Q2.} Si $x \in [20\,;40]$ alors $15 \leq 0{,}75x \leq 30$, donc $20 \leq f(x) \leq 35$, et en particulier $f(x) \in [20\,;40]$. Recurrence : $s_0 = 40 \in [20\,;40]$, et la stabilite conclut.

\textbf{Q3.} $f(x) - x = 5 - 0{,}25x = -\dfrac{x-20}{4}$, negatif sur $[20\,;40]$. Donc $s_{n+1} - s_n \leq 0$ : la suite est decroissante. Minoree par $20$, elle converge.

\textbf{Q4.} $f$ est continue (affine) et la limite $\ell$ verifie $0{,}75\ell + 5 = \ell$, soit $0{,}25\ell = 5$ et $\ell = 20$. La concentration tend vers celle de l'eau ajoutee : le bassin se rapproche indefiniment de $20$ g$\cdot$L$^{-1}$ sans jamais descendre en dessous, ce qui est coherent puisqu'on n'ajoute jamais d'eau moins salee que $20$.

\textbf{Q5.} $e_{n+1} = s_{n+1} - 20 = 0{,}75\,s_n + 5 - 20 = 0{,}75(s_n - 20) = 0{,}75\,e_n$. La suite $(e_n)$ est geometrique de raison $0{,}75$ et de premier terme $e_0 = 20$, donc $s_n = 20 + 20 \times 0{,}75^{\,n}$.

On cherche $s_n < 21$, soit $20 \times 0{,}75^n < 1$, c'est-a-dire $0{,}75^n < 0{,}05$. Or $0{,}75^{10} \approx 0{,}0563 > 0{,}05$ et $0{,}75^{11} \approx 0{,}0422 < 0{,}05$ : la concentration passe sous $21$ g$\cdot$L$^{-1}$ au bout de \textbf{11} jours.
\end{corrige}
